\documentclass[10pt,a4paper]{amsart}
\usepackage[margin=19mm]{geometry}
\usepackage{amsmath,amssymb,mathtools}
\usepackage[scr=rsfs,cal=euler]{mathalfa}
\usepackage{xfrac}
\usepackage[unicode,hidelinks,bookmarks=false]{hyperref}
\newcommand{\ie}{i.e.\ }
\newcommand{\cf}{cf.\ }
\newcommand{\resp}{respectively\ }
\newcommand{\Subsection}{\subsection}
\providecommand{\nobreakdash}{\nobreak\hbox{-}}
\setlength{\emergencystretch}{2em}
\multlinegap0pt
\usepackage{amssymb}
\usepackage{tikz}
\usepackage[normalem]{ulem}
\usepackage{xcolor}
\newtheorem{lemma}{Lemma}
\title{Approximation theorems in bilipschitz invariant theory}
\author{Jameson Cahill, Joseph W. Iverson, Dustin G. Mixon, Nathan Willey}
\date{Source: arXiv:2603.23643v1 (CC BY 4.0)}
\begin{document}
\maketitle

\noindent\emph{Source and licence.} Approximation theorems in bilipschitz invariant theory. Version arXiv:2603.23643v1 (CC BY 4.0), distributed by arXiv under the Creative Commons Attribution 4.0 International licence, http://creativecommons.org/licenses/by/4.0/, as recorded in the arXiv licence metadata for this exact version and shown on the abstract page https://arxiv.org/abs/2603.23643v1. Full licence notice: \texttt{licenses/arxiv-2603.23643-CC-BY-4.0.txt}.\par\smallskip

\noindent\textit{Editorial scope.} Counted unit: the article's lemma lem.ksl (source0.tex lines 1223--1232). The article's complete original proof of it is delivered with it (source0.tex lines 1330--1394, one \texttt{proof} environment of 65 lines, which the article places away from the statement). No mathematics is written, repaired or re-derived: the statement and the proof are the article's own lines, in the article's own notation, and no retained line is substituted or re-flowed.  Retained uncounted as the source-local context the proof needs: lines 1127--1131.  Nothing was omitted from any retained span, so the package carries no omission record.  No retained line contains a citation, so the wrapper carries no bibliography and no bibliography entry.  No retained line contains a figure environment, an image-inclusion command or drawing code, so no image is delivered and the package declares no asset dependency.  Editorial formatting only, in addition: an AMS \texttt{amsart} preamble that restates the article's own declarations and macros used by the retained lines, namely the environments \texttt{lemma} on separate counters, together with the standard packages those lines need.  The retained environments print in this document as Lemma 1 in order of appearance, whereas the article prints the same items with the numbering of its own full document.  The complete native source of the article as distributed in the arXiv e-print for this exact version is bundled unchanged as \texttt{sources/197093/source0.tex}.
\begin{equation}
\label{eq.stratification of PL}
\operatorname{PL}_{\operatorname{h}}
=\bigcup_\mathcal{H}\operatorname{PL}_{\operatorname{h}}(\mathcal{H}),
\end{equation}

\begin{lemma}
\label{lem.ksl}
Take any $d>2$ and put $G := \{ \pm  I \} \leq \operatorname{O}(d)$.
Given $f\in\operatorname{PL}_{\operatorname{h}}^G$, denote
\[
K(f)
:=\{\,x \in \mathbb{R}^d : \text{$f$ is not differentiable at $x$}\,\}.
\]
If $K(f)$ is not a union of finitely many hyperplanes, then $f \notin \overline{\operatorname{span}\mathcal{M}}$.
\end{lemma}

\begin{proof}[Proof of Lemma~\ref{lem.ksl}]
First, we prove an intermediate claim, which is a special case of the contrapositive:
For every nonzero $f\in\operatorname{span}\mathcal{M}$, it holds that $K(f)$ is a union of finitely many hyperplanes.
Take any nonzero $f\in\operatorname{span}\mathcal{M}$, and consider any decomposition $f=\sum_{i=1}^n c_i|\langle\cdot,y_i\rangle|$ into max filters $\{|\langle\cdot,y_i\rangle|\}_{i=1}^n$ that are linearly independent as functions.
We claim that $K(f)=\bigcup_{i=1}^n\operatorname{span}\{y_i\}^\perp$.
For notational convenience, we use $U$ to denote the right-hand set.
The $\subseteq$ direction follows from the fact that any point of nondifferentiability in $f$ is a point of nondifferentiability for some max filter $|\langle\cdot,y_i\rangle|$.
For the $\supseteq$ direction, consider any $x\in U$.
Put $\delta_i:=\operatorname{sign}\langle x,y_i\rangle\in\{-1,0,1\}$.
Select a generic $v\in\mathbb{R}^d$.
Then for all $t\in\mathbb{R}$ with sufficiently small absolute value, it holds that
\begin{align*}
f(x+tv) 
= \sum_{i=1}^n c_i|\langle x+tv,y_i\rangle|
&= \sum_{\substack{i\in[n]\\[1 pt]y_i\not\perp x}} c_i|\langle x+tv,\delta_iy_i\rangle|+\sum_{\substack{i\in [n]\\[1 pt]y_i\perp x}} c_i|\langle x+tv,y_i\rangle|\\
&= \sum_{\substack{i\in[n]\\[1 pt]y_i\not\perp x}} c_i\langle x+tv,\delta_iy_i\rangle+\sum_{\substack{i\in[n]\\[1 pt]y_i\perp x}} c_i|\langle tv,y_i\rangle|\\
&=f(x)+t\sum_{\substack{i\in[n]\\[1 pt]y_i\not\perp x}} c_i\langle v,\delta_iy_i\rangle+|t|\sum_{\substack{i\in[n]\\[1 pt]y_i\perp x}} c_i|\langle v,y_i\rangle|.
\end{align*}
Since $x\in U$, there necessarily exists at least one $i$ for which $y_i\perp x$.
Furthermore, the coefficient of $|t|$ is nonzero since $v$ is generic.
It follows that $f$ is not differentiable at $x$ in the direction of $v$, and so $x\in K(f)$, as claimed.

We are now ready to prove the lemma.
Take any $f\in\operatorname{PL}_{\operatorname{h}}^G$ for which $K(f)$ is not a union of finitely many hyperplanes.
Since $f\in\operatorname{PL}_{\operatorname{h}}$, the stratification~\eqref{eq.stratification of PL} gives that there exists a set $\mathcal{H}$ of hyperplanes for which $f\in\operatorname{PL}_{\operatorname{h}}(\mathcal{H})$.
For each $H\in\mathcal{H}$, select a unit vector $z_H\in H^\perp$, and let $T\leq\operatorname{PL}_{\operatorname{h}}(\mathcal{H})$ denote the span of the resulting max filters $\{|\langle\cdot,z_H\rangle|\}_{H\in\mathcal{H}}$.
Note that for every linear functional $\ell\in(\mathbb{R}^d)^*$, the function $f+\ell\in\operatorname{PL}_{\operatorname{h}}(\mathcal{H})$ is nonzero with $K(f+\ell)=K(f)$.
Meanwhile, by our intermediate claim, every nonzero member of $T$ fails to be differentiable over an entire linear hyperplane in $\mathcal{H}$.
Since $f+(\mathbb{R}^d)^*$ avoids $T$, and since $\kappa$ defines a norm on $\operatorname{PL}_{\operatorname{h}}(\mathcal{H})/(\mathbb{R}^d)^*$, it follows that
\[
c
:=\inf_{g\in T}\kappa(g-f)
>0.
\]
Next, we show that for every linear combination $\sum_{i=1}^n c_ig_i$ of max filters $\{g_i\}_{i=1}^n$,
\begin{equation}
\label{eq.claimed bound}
\bigg\|\sum_{i=1}^n c_ig_i - f\bigg\|_{\operatorname{Lip}}
\geq\frac{c}{2},
\end{equation}
so that $f$ does not reside in the Lipschitz closure of the span of max filters.
Let $I\subseteq[n]$ denote the indices for which $g_i\in T$, and put
$g:=\sum_{i\in I} c_ig_i$
and $h:=\sum_{i\in[n]\setminus I} c_ig_i$.
Then $g-f\in\operatorname{PL}_{\operatorname{h}}(\mathcal{H})$ and $\kappa(g-f)\geq c$.
Furthermore, since $g_i\not\in T$ for every $i\in[n]\setminus I$, none of the linear hyperplanes of non-differentiability in $h$ coincide with any of the linear hyperplanes in $\mathcal{H}$.
It follows that we may select $x,y\in\mathbb{R}^d$ from the interiors of adjacent chambers in $\mathcal{T}(\mathcal{H})$ in such a way that
\[
\|\nabla (g-f)(x)-\nabla (g-f)(y)\|=\kappa(g-f)
\qquad
\text{and}
\qquad
\nabla h(x)=\nabla h(y).
\]
This then gives
\begin{align*}
c
\leq\kappa(g-f)
=\|\nabla (g-f)(x)-\nabla (g-f)(y)\|
&=\|\nabla h(x)-\nabla h(y) + \nabla (g-f)(x)-\nabla (g-f)(y)\|\\
&\leq\|\nabla(h+g-f)(x)\|+\|\nabla(h+g-f)(y)\|\\
&\leq 2\|h+g-f\|_{\operatorname{Lip}},
\end{align*}
which rearranges to the claimed bound~\eqref{eq.claimed bound}.
\end{proof}

\end{document}
